\documentclass[11pt,a4paper]{article}
\usepackage[a4paper,margin=2.25cm,headheight=15pt]{geometry}
\usepackage{fontspec}
\setmainfont{Arial}
\usepackage{polyglossia}
\setdefaultlanguage{brazil}
\usepackage{xcolor}
\definecolor{azul}{HTML}{174A6E}
\definecolor{cinza}{HTML}{4B5563}
\usepackage{titlesec}
\titleformat{\section}{\Large\bfseries\color{azul}}{\thesection}{0.6em}{}
\titleformat{\subsection}{\large\bfseries\color{azul}}{\thesubsection}{0.6em}{}
\titleformat{\subsubsection}{\normalsize\bfseries\color{azul}}{\thesubsubsection}{0.6em}{}
\setcounter{secnumdepth}{0}
\usepackage{enumitem}
\setlist{itemsep=2pt,topsep=4pt,leftmargin=1.6em}
\usepackage{microtype}
\usepackage{hyperref}
\usepackage{tabularx}
\hypersetup{colorlinks=true,linkcolor=azul,urlcolor=azul,pdftitle={Guia de leitura do projeto CHM-SIBGRAPI}}
\usepackage{fancyhdr}
\pagestyle{fancy}\fancyhf{}
\fancyhead[L]{\small\color{cinza}CHM-SIBGRAPI | Guia de leitura}
\fancyhead[R]{\small\color{cinza}Redes neurais e sensoriamento remoto}
\fancyfoot[C]{\thepage}
\setlength{\parskip}{5pt}\setlength{\parindent}{0pt}
\begin{document}
\section{Guia de leitura do projeto CHM-SIBGRAPI}
Este documento explica, em linguagem introdutória, o que o projeto faz, como os dados viram previsões e para que servem os principais scripts e pastas. Ele foi escrito a partir do código, dos READMEs, da configuração do pipeline e dos metadados do conjunto de dados que estão neste repositório.

\begin{quote}\color{cinza}
\textbf{Em uma frase:} o projeto combina imagens de radar e imagens ópticas de satélite para estimar a altura da vegetação em uma floresta da Guiana Francesa, usando medições LiDAR como referência para ensinar e avaliar os modelos.\par
\end{quote}
\subsection{1. O problema que o projeto tenta resolver}
Uma floresta tem propriedades que não são fáceis de medir em grandes áreas. Aqui, o alvo principal é o \textbf{CHM} (\emph{Canopy Height Model}): um mapa de altura do topo da vegetação, em metros. O CHM de referência é construído a partir de dados LiDAR: em termos simples, subtrai-se o modelo do terreno (DTM) do modelo da superfície (DSM). O LiDAR fornece a resposta conhecida durante o treinamento; o modelo tenta aprender a estimá-la usando sensores que cobrem áreas maiores.

O projeto também transforma altura em \textbf{AGB} (\emph{Above-Ground Biomass}, biomassa acima do solo). Essa conversão usa uma relação alométrica, isto é, uma equação empírica que relaciona altura da vegetação a biomassa. AGB é, portanto, uma estimativa derivada do CHM e da equação adotada; não é uma medição direta feita pela rede neural.

\subsection{2. Sensoriamento remoto: os dados de entrada}
\subsubsection{Radar SAR}
\textbf{SAR} significa radar de abertura sintética. O sensor envia micro-ondas e registra o eco que retorna. Diferentemente de uma câmera comum, o radar mede uma resposta eletromagnética que depende da estrutura, umidade, orientação e composição dos alvos. Ele pode observar através de nuvens, uma vantagem importante em regiões tropicais.

O projeto usa dados \textbf{PolInSAR} em duas bandas, \textbf{L} e \textbf{P}, para duas aquisições (05/06 e 06/03, identificadas no código como \url{tropi0506} e \url{tropi0603}). Quatro polarizações aparecem como \url{HH}, \url{HV}, \url{VH} e \url{VV}: as letras indicam a orientação da onda transmitida e recebida. O processamento forma uma \textbf{coerência complexa normalizada} para cada polarização. Pense nela como uma medida de semelhança entre as respostas do radar nas duas aquisições; por ser complexa, ela contém parte real e imaginária, que carregam informação de amplitude e fase.

O radar não vira simplesmente quatro números. O modelo recebe pequenas matrizes espaciais (\emph{patches}), preservando padrões ao redor de cada local. A geometria dessas imagens é especialmente importante: o \textbf{back-geocoding} usa grades de localização chamadas \url{.grille} para trazer as imagens SAR para uma grade comum com as demais fontes. A política atual exige interseção geométrica válida das quatro grades, evitando usar pixels extrapolados ou fora das imagens.

\subsubsection{Landsat (óptico)}
Landsat é um satélite que mede luz refletida em diferentes faixas do espectro. O conjunto usa seis bandas de reflectância e dois índices calculados a partir delas: \textbf{NDVI} e \textbf{EVI}, indicadores relacionados à vegetação. Nuvens, saturação e pixels sem dados podem invalidar uma observação, então o pipeline usa as máscaras de qualidade (\url{QA_PIXEL} e, quando disponível, \url{QA_RADSAT}).

O mosaico atual seleciona, para cada célula de 30 m, uma cena considerada válida e de melhor qualidade. As bandas de uma célula vêm da mesma cena, para não misturar datas dentro daquele pixel. O mosaico continua sendo \textbf{multitemporal}: diferentes células podem vir de datas diferentes. O arquivo \url{metadata.json} registra as cenas e a data selecionada. Essas imagens não são necessariamente simultâneas às aquisições TropiSAR.

\subsubsection{LiDAR e CHM}
LiDAR mede distâncias usando pulsos de laser e permite construir superfícies de terreno e de copa com bastante detalhe. O projeto valida sistema de coordenadas, resolução, valores inválidos e alturas extremas antes de usar esses dados como referência. O CHM é o \textbf{alvo} que o modelo tenta prever; não deve ser confundido com uma entrada óptica ou de radar.

\subsection{3. O que é uma rede neural, sem pressupor conhecimento prévio}
Uma rede neural é uma função com muitos parâmetros ajustáveis. Ela recebe números (aqui, matrizes de bandas e canais de sensores), produz uma previsão e compara essa previsão com o valor de referência. No treinamento, um algoritmo ajusta os parâmetros para diminuir os erros. Repetir esse processo em muitos exemplos permite que o modelo aprenda relações úteis nos dados.

Uma \textbf{época} é uma passagem pelo conjunto de treinamento. Um \textbf{batch} é um pequeno grupo de exemplos processado de uma vez. A \textbf{loss} (função de perda) resume o erro que o treinamento tenta reduzir. O conjunto de \textbf{treino} ajusta os parâmetros; o de \textbf{validação} ajuda a escolher quando parar e qual versão do modelo guardar. Ele não deve ser confundido com um teste final independente.

O projeto padroniza as entradas usando estatísticas do conjunto de treino. Isso coloca canais com escalas diferentes em faixas mais comparáveis. A padronização do alvo CHM também é usada internamente, mas as previsões são reconvertidas para metros.

\subsubsection{Arquitetura principal: TCCT}
O modelo principal é identificado como \textbf{TCCT}. É uma arquitetura de fusão: cada tipo de sensor tem um caminho de processamento, e depois as informações são combinadas para prever o CHM.

\begin{itemize}
\item \textbf{CVNN / convoluções complexas:} o caminho do radar trata números complexos (partes real e imaginária) em vez de descartar diretamente a fase. Convoluções procuram padrões locais, como formas e texturas, em cada patch.
\item \textbf{Encoder óptico:} extrai padrões das bandas Landsat e dos índices de vegetação.
\item \textbf{Coatenção trimodal:} mecanismo que permite aos caminhos de radar L, radar P e óptico ponderar informações uns dos outros. A intenção é aprender quais fontes ajudam mais em cada exemplo.
\item \textbf{Cabeça global/local:} combina um resumo do patch com detalhes espaciais para produzir um mapa CHM previsto dentro do patch.
\item \textbf{Gating modal:} pesos/diagnósticos associados à contribuição das modalidades são exportados para inspeção. Eles ajudam a entender o comportamento do modelo, mas não provam causalidade nem explicam sozinhos a decisão.
\end{itemize}
O código contém variantes para comparação: TCCT sem coatenção, uma CNN com fusão antecipada e um ViT denso (\emph{Vision Transformer}). Também existe uma baseline \textbf{Random Forest}, que é um conjunto de árvores de decisão e não uma rede neural.

\subsubsection{Função de perda e duas fases}
A loss do TCCT combina mais de um critério. Ela considera erro nos pixels de CHM, erro no resumo por patch, diferenças de gradiente/estrutura e concordância estatística entre valores previstos e observados. Em áreas de vegetação alta, o código pode dar peso adicional aos exemplos acima do limiar configurado. Isso tenta equilibrar diferentes aspectos do problema, mas não significa que o modelo esteja correto só porque a loss caiu.

O treinamento do TCCT pode ter duas fases:

\begin{enumerate}
\item \textbf{Fase A supervisionada:} aprende a prever o CHM usando as referências LiDAR e a loss de imagem.
\item \textbf{Fase B com termo RVoG:} continua o ajuste acrescentando uma restrição inspirada na física do radar, quando há dados geométricos válidos.
\end{enumerate}
\textbf{RVoG} (\emph{Random Volume over Ground}) é um modelo físico simplificado para descrever como o sinal de radar interage com vegetação e solo. Nesta implementação, geometria como \url{kz} (número de onda vertical) e ângulo de incidência ajuda a calcular um termo de consistência. Esse termo é adicionado gradualmente na Fase B. A ideia é combinar aprendizado a partir de exemplos com uma restrição física plausível; não é uma medição direta da altura.

\subsection{4. Fluxo do pipeline principal}
A ordem conceitual do arquivo \url{Pipeline_TCCT...RETOMAVEL.py} é:

\begin{enumerate}
\item \textbf{Configuração e ambiente:} escolhe diretórios, bibliotecas, dispositivo disponível (GPU CUDA, Apple MPS ou CPU), parâmetros, sementes e caminhos do projeto.
\item \textbf{Auditoria dos arquivos:} inventaria dados e metadados, valida formatos, sistema de coordenadas e layouts binários. O objetivo é encontrar incompatibilidades antes das etapas pesadas.
\item \textbf{Preparação LiDAR/CHM:} valida o DTM e DSM, produz o CHM canônico e relatórios/figuras diagnósticas.
\item \textbf{Back-geocoding SAR:} mapeia as imagens SAR para a grade geográfica comum usando as \url{.grille}; registra mapas de linha/coluna, máscaras e auditoria geométrica.
\item \textbf{Landsat e coerência:} cria o mosaico óptico controlado por qualidade e calcula mapas de coerência normalizada do radar.
\item \textbf{Dataset:} junta entradas SAR, Landsat, CHM e metadados em patches. Cada exemplo contém matrizes de entrada, CHM de referência, máscara de pixels válidos e informações de qualidade/geometria. No dataset v6 atual há \textbf{10.500 amostras}.
\item \textbf{Geometria RVoG:} prepara ou audita parâmetros físicos para a banda P. A configuração atual indica \url{ha_p_real}, não permite calibrar RVoG livremente a partir do conjunto de treino, mas permite calibrar atenuação/signo com \url{kz} real.
\item \textbf{Folds espaciais:} divide os exemplos por blocos/regiões, e não por patches sorteados sem considerar localização. Isso reduz a chance de treino e validação conterem vizinhos quase idênticos. O projeto usa 5 folds e audita também o efeito de um buffer espacial.
\item \textbf{Treino e avaliação:} calcula estatísticas somente no treino, cria loaders, treina os modelos, salva checkpoints e avalia previsões.
\item \textbf{Experimentos multiseed:} repete as combinações com sementes 42, 314 e 2026. Isso mostra quanto o resultado pode variar por inicialização aleatória. Cinco folds × três seeds dão 15 execuções por configuração.
\item \textbf{Comparações e produtos:} reúne resultados, executa baselines, salva previsões fora da amostra, gera tabelas, gráficos e mosaicos CHM/AGB.
\item \textbf{Retomada:} checkpoints e estados por época permitem retomar treinamentos interrompidos. Os resumos e arquivos \url{estado_execucao*.json} mostram o andamento.
\end{enumerate}
A configuração salva em \url{Resultados/TCCT_v3_2_SIBGRAPI/config_pipeline.json} é uma fotografia dos parâmetros utilizados. Entre eles: patches de 30 m, janela de coerência 5, batch size 24, 5 folds, três seeds e limites de épocas para as duas fases. Os valores são escolhas experimentais, não leis gerais.

\subsection{5. Validação espacial: por que importa}
Em imagens geográficas, dois patches vizinhos costumam se parecer muito. Se um vai para treino e outro para validação, o modelo pode parecer bom por ter visto praticamente o mesmo local. Isso é \textbf{vazamento espacial}: a avaliação fica otimista e não representa bem a previsão em uma região nova.

Por isso, o pipeline agrupa por blocos espaciais e faz validação cruzada em cinco folds. Em cada fold, uma região fica para validação e as demais para treino. O treinamento multiseed repete o processo com diferentes sementes. Os manifests \url{manifesto_fold_*.csv} registram quais amostras pertencem a cada divisão; seus hashes ajudam a detectar mudanças.

O guia de retomada também menciona um experimento \textbf{holdout 70/20/10}: aproximadamente 70\% treino, 20\% validação e 10\% teste final intocado, por faixas de blocos. É uma análise adicional, recomendada depois de congelar os hiperparâmetros. Não substitui os folds espaciais principais.

\subsection{6. Como ler as métricas}
As principais métricas de regressão aparecem nos resumos:

\begin{itemize}
\item \textbf{MAE:} erro absoluto médio, em metros para o CHM. Menor é melhor e costuma ser fácil de interpretar.
\item \textbf{RMSE:} raiz do erro quadrático médio, também em metros. Penaliza mais erros grandes.
\item \textbf{R²:} mede quanto da variação dos valores é explicada em relação a uma referência média. Quanto mais perto de 1, melhor; pode ser negativo se o modelo for pior que prever a média.
\item \textbf{CCC:} coeficiente de concordância, que considera correlação e também diferenças de média/escala. Mais perto de 1 indica melhor concordância.
\item \textbf{Média ± desvio padrão entre folds/seeds:} resume desempenho e variabilidade. Uma média boa com grande desvio sugere instabilidade entre regiões ou inicializações.
\end{itemize}
Não escolha um modelo olhando apenas uma métrica. Compare desempenho, variabilidade, distribuição dos erros e a forma como a validação foi feita. Uma métrica alta não elimina problemas de qualidade dos dados ou de representatividade geográfica.

\textbf{OOF} (\emph{out-of-fold}) significa que cada previsão é produzida por um modelo que não treinou naquela amostra/região. Essas previsões podem ser reunidas em mosaicos e comparadas de forma mais justa. As figuras pareadas e os scripts de mosaico usam predições salvas; eles não precisam treinar o modelo novamente.

\subsection{7. Mapa rápido dos arquivos de código}
Pasta \href{Codigo}{Codigo}:

\begin{itemize}
\item \url{Pipeline_TCCT_v3_2_1_SIBGRAPI_REV_GRILLE3D_RETOMAVEL.py}: núcleo do projeto. Faz auditorias e processamento, monta dataset, define modelos, loss, treinamento, folds, avaliação e produtos derivados.
\item \url{Pipeline_TCCT_v3_2_1_SIBGRAPI_REV_GRILLE3D_PORTATIL_RETOMAVEL.ipynb}: interface de notebook para executar etapas do pipeline. O notebook portátil e o \url{.py} correspondem ao mesmo fluxo geral; para repetir a execução documentada, siga o README de retomada.
\item \url{GUIA_CIENTIFICO_TCCT_v3_2_1.ipynb}: material científico de apoio sobre método e decisões experimentais.
\item \url{Calcular_AGB_Todos_Modelos_Predicoes_Salvas.py}: carrega predições já salvas, alinha amostras de validação com células do inventário florestal, calcula métricas CHM/AGB e gera tabelas/figuras. Não é o treinamento central.
\item \url{TCCT_AGB_Allometry_SupportMatched.py}: calibra/avalia uma relação alométrica com suporte espacial compatível entre CHM e células do inventário. A preocupação com suporte evita comparar uma média de área com uma medida de área diferente sem ajuste.
\item \url{TCCT_Gerar_Mosaico_OOF_CHM_e_AGB.py}: reconstrói mosaicos de previsões OOF a partir dos checkpoints/folds e pode avaliar biomassa no suporte do inventário.
\item \url{Gerar_Figuras_Pareadas_Predicoes_Salvas.py}: valida o alinhamento das predições por amostra e gera gráficos comparativos usando resultados existentes, sem retreinar.
\item \url{Coletar_Artefatos_Multiseed_TCCT.py}: coleta e audita resultados de múltiplas seeds/folds, com hashes e opção de pacote ZIP.
\item \url{Preparar_Pacote_Reprodutibilidade_TCCT_macOS_corrigido.py}: inventaria arquivos, confere a completude das execuções e prepara artefatos de reprodutibilidade. O README correspondente explica modos de cópia; não é código de treinamento.
\item \url{CHANGELOG_TCCT_v3_2_1.md}: alterações da versão.
\item \url{MANIFESTO_SHA256_TCCT_v3_2_1.txt}: hashes de arquivos para conferir integridade.
\end{itemize}
Os notebooks podem conter cópias ou chamadas das etapas do pipeline. O arquivo \url{.py} é a referência mais fácil para localizar definições e funções; os notebooks são uma maneira interativa de executar e documentar o processo. O README de retomada descreve o uso operacional do Passo 20 e do holdout.

\subsection{8. Como interpretar as pastas de dados e resultados}
\begin{itemize}
\item \url{Dados_Brutos/}: fontes originais. São entradas, não resultados do modelo.
\item \url{Dados_Processados/}: CHM, mosaicos, mapas de coerência, back-geocoding e geometria derivados.
\item \url{Dataset/TCCT_dataset_v6_grille3d/}: índice e metadados do dataset usado. \url{metadata.json} documenta a versão, número de amostras e decisões de qualidade; \url{index.csv} cataloga patches.
\item \url{Auditoria/}: relatórios sobre arquivos, geometrias, filtros, coordenadas e reprodutibilidade.
\item \url{Resultados/Metricas/}: resultados agregados e detalhados por modelo/fold, estados de execução e resumos.
\item \url{Resultados/TCCT_v3_2_SIBGRAPI/}: checkpoints, históricos, configurações e resultados do pipeline. \url{multiseed/} organiza execuções por seed/fold.
\item \url{Resultados/Figuras/}: diagnósticos, comparações e visualizações geradas.
\item Arquivos \url{.pth}: pesos/checkpoints de redes neurais. Eles guardam parâmetros aprendidos; não são código-fonte nem imagens legíveis diretamente.
\item Arquivos \url{.csv}: tabelas; \url{.json}: configurações e relatórios; \url{.npy}/\url{.npz}: matrizes numéricas; \url{.tif}: raster geográfico.
\end{itemize}
Os arquivos existentes são resultados de execuções específicas; o fato de estarem na pasta não garante que todas as combinações planejadas terminaram. Consulte os resumos e estados de execução antes de afirmar que o experimento está completo.

\subsection{9. Um glossário curto}
\begin{center}\small\begin{tabularx}{\textwidth}{|X|X|}\hline
Termo & Em palavras simples \\ \hline
Patch & Pequeno recorte de imagem usado como um exemplo do modelo. \\ \hline
Banda & Faixa do espectro medida por um sensor. \\ \hline
Polarização & Orientação da onda do radar ao transmitir/receber. \\ \hline
Coerência & Semelhança complexa entre respostas de radar de aquisições relacionadas. \\ \hline
Back-geocoding & Mapeamento dos pixels SAR para suas posições geográficas. \\ \hline
Tensor & Matriz numérica que o modelo processa; pode ter canais, linhas e colunas. \\ \hline
Parâmetro/peso & Número ajustável dentro da rede, aprendido durante o treinamento. \\ \hline
Seed/semente & Valor que controla aleatoriedade para tornar uma execução repetível. \\ \hline
Fold & Uma das divisões da validação cruzada. \\ \hline
Checkpoint & Arquivo com estado/pesos do modelo para avaliar ou retomar treinamento. \\ \hline
Ablation/ablação & Experimento que remove uma parte do modelo para medir sua contribuição. \\ \hline
Baseline & Método de referência simples ou alternativo para comparação. \\ \hline
Vazamento & Informação de validação/teste que influencia o treinamento ou deixa a avaliação artificialmente fácil. \\ \hline
CRS & Sistema de referência de coordenadas que informa onde os pixels ficam na Terra. \\ \hline
Raster & Imagem em grade geográfica, com um valor por pixel. \\ \hline
\end{tabularx}\end{center}
\subsection{10. Um roteiro amigável para explorar o código}
\begin{enumerate}
\item Leia este guia e o \href{README_TCCT_v3_2_1_RETOMADA.md}{README de instalação e retomada}.
\item Abra \url{Resultados/TCCT_v3_2_SIBGRAPI/config_pipeline.json} para ver os parâmetros e \url{Dataset/TCCT_dataset_v6_grille3d/metadata.json} para entender de onde vieram as 10.500 amostras.
\item No pipeline \url{.py}, localize os títulos de seção começando por \url{0.} até \url{9.}. Leia primeiro as seções 4 a 7: dataset, folds, arquitetura, loss e treinamento.
\item Depois leia as etapas de preparação (1 a 3B) para acompanhar como as fontes são alinhadas.
\item Por fim, examine um \url{resumo_final_*.csv}, um \url{resumo_por_fold_*.csv} e os históricos. O primeiro resume; o segundo revela variação geográfica; o histórico mostra o treinamento época a época.
\item Ao comparar modelos, procure resultados por fold e seed, e confira se as previsões estão alinhadas às mesmas amostras.
\end{enumerate}
\subsubsection{Limites desta explicação}
Este guia explica a intenção e a estrutura do código, mas não substitui uma revisão científica dos pressupostos, das escolhas de hiperparâmetros ou da validade externa do estudo. A qualidade do modelo depende da qualidade/alinhamento dos dados e de quão bem as regiões de validação representam locais onde se deseja aplicar o método.

\end{document}
